% Part:sets-functions-relations
% Chapter: sets
% Section: schroder-bernstein

\documentclass[../../../include/open-logic-section]{subfiles}

\begin{document}

\olfileid{sfr}{siz}{sb}

\olsection{કદનો ખ્યાલ અને શ્રેડર--બર્નસ્ટાઇન પ્રમેય}

\begin{explain}
આ એક સહજ વિચાર છે: જો $A$ એ $B$ કરતાં મોટો ન હોય અને $B$ એ
$A$ કરતાં મોટો ન હોય, તો $A$ અને $B$ સામ્ય ગણો છે. સાચું કહીએ તો,
જો આ વિચાર \emph{ખોટો} હોત, તો ગણો વચ્ચેના “કદ”ની તુલના સાથે
ગણસામ્યના આપણા વ્યાખ્યાયિત ખ્યાલનો કોઈ સંબંધ છે તે વિચારને ભાગ્યે જ
યોગ્ય ઠરાવી શકાત!{} સદભાગ્યે આ સહજ વિચાર સાચો છે. શ્રેડર--બર્નસ્ટાઇન
પ્રમેય તેને યોગ્ય ઠરાવે છે.
\end{explain}

\begin{thm}[શ્રેડર--બર્નસ્ટાઇન]
	\ollabel{thm:schroder-bernstein}
	જો $\cardle{A}{B}$ અને $\cardle{B}{A}$ હોય,
	તો $\cardeq{A}{B}$.
\end{thm}

\begin{explain}
બીજા શબ્દોમાં, જો $A$માંથી~$B$માં કોઈ !!a{injection}, અને
$B$માંથી~$A$માં કોઈ !!a{injection} હોય, તો $A$માંથી~$B$માં
કોઈ !!a{bijection} હોય છે.

પણ આ પરિણામની સાબિતી ખરેખર ઘણી \emph{મુશ્કેલ} છે. કૅન્ટૉરે
પરિણામ જણાવ્યું હતું, પરંતુ બીજાઓએ તે સાબિત કર્યું હતું.
\footnote{ઇતિહાસ વિશે વધુ જાણવા, દાખલા તરીકે,
\citet[pp.~165--6]{Potter2004} જુઓ.}
\oliflabeldef{sfr:cardinals:card-sb:sec}{શ્રેડર--બર્નસ્ટાઇનને
\emph{સાબિત} કરવાની સ્થિતિમાં આપણે
\olref[sfr][cardinals][card-sb]{sec}માં જ આવીશું.}{}%
અત્યારે તમારે તેને વિશ્વાસપૂર્વક સ્વીકારવું (અને સ્વીકારવું જ) પડશે.

સદભાગ્યે, શ્રેડર--બર્નસ્ટાઇન \emph{સાચું} છે અને તે આપણે વ્યાખ્યાયિત
કરેલા સંબંધો, એટલે કે $\cardeq{A}{B}$ અને $\cardle{A}{B}$ને “કદ”
સાથે કંઈક સંબંધ ધરાવતા માનીએ છીએ તેને સમર્થન આપે છે. વળી,
શ્રેડર--બર્નસ્ટાઇન ઘણું \emph{ઉપયોગી} છે. બે સામ્ય ગણો વચ્ચેનું
!!a{bijection} વિચારવું મુશ્કેલ હોઈ શકે છે. શ્રેડર--બર્નસ્ટાઇન પ્રમેય
આપણને તુલનાને બે ભાગમાં વહેંચવા દે છે, જેથી આપણે માત્ર પહેલા ગણમાંથી
બીજામાં જતું અને પછી બીજામાંથી પહેલામાં જતું !!a{injection} વિચારવું પડે.
\end{explain}
\end{document}
